
Pour cela, on considère le complexe :
\[
0\longrightarrow \wedge^r(\underline{E})\otimes_{\underline{O}_S}\mathbb S(\underline{E})\longrightarrow
\wedge^{r-1}(\underline{E})\otimes_{\underline{O}_S}\mathbb S(\underline{E})\longrightarrow\ldots\longrightarrow
\underline{E}\otimes_{\underline{O}_S}\mathbb S(\underline{E})\longrightarrow \mathbb S(\underline{E})\longrightarrow 0
\]
où $\mathbb S(\underline{E})$ est l'algèbre symétrique de $\underline{E}$. C'est une résolution de $\underline{O}_S$ sur $\mathbb S(E)$ (VI 1.11). Le complexe associé sur $X$, noté $\underline{L}_\bullet$, a pour $k$-ième terme $\wedge^k(f^*(E))$ et est donc une résolution de $\underline{O}_S$ par des $\underline{O}_X$-Modules localement libres de type fini ; on a donc \(\underline{L}_\bullet\otimes_{\underline{O}_X}\underline{F}=\underline{O}_S\otimes^{\mathbb L}_{\underline{O}_X}\underline{F}\). Par suite :
\[
\begin{aligned}
g^*(g_*(F))
&=\cl(\underline{L}_\bullet\otimes_{\underline{O}_X}\underline{F})
=\sum_{i=0}^r(-1)^i\cl(\underline{L}_i\otimes_{\underline{O}_X}\underline{F})\\
&=\sum_{i=0}^r(-1)^i\cl(\wedge^i(\underline{E})\otimes_{\underline{O}_S}\underline{F})
=F\cdot\lambda_{-1}(E).
\end{aligned}
\]
La proposition en résulte par additivité.

\subsection*{2. Fibrés principaux sous les tores déployés}

2.1. Soient $S$ un schéma noethérien, et $T$ un tore sur $S$. On suppose $T$ déployé; on a donc un isomorphisme $T\simeq D_S(M)$, où $M$ est un groupe abélien libre de rang $r$, soit encore $T\simeq \underline{G}_{m\,S}^r$. Soit $P$ un torseur sous $T$. Alors $P$ est affine sur $S$, défini par une $\underline{O}_S$-Algèbre $\underline{A}$ graduée de type $M$ ; si on pose $\underline{A}=\bigoplus_{m\in M}\underline{A}_m$, alors $\underline{A}_m$ est un $\underline{O}_S$-Module inversible pour tout $m$, et, pour $m,m'\in M$, l'homomorphisme $\underline{A}_m\otimes_{\underline{O}_S}\underline{A}_{m'}\longrightarrow \underline{A}_{m+m'}$ défini par la multiplication de $\underline{A}$ est un isomorphisme (SGA 3 VIII 4.1). On définit donc un homomorphisme multiplicatif $\eta:M\longrightarrow K^{\bullet}(S)$ en associant à $m\in M$ la classe de $\underline{A}_m$ dans $K^{\bullet}(S)$.

